\section{Runtime Typing Consequences in Singular Operations}

At this point the basic pieces are in place: names are bindings, singular values are heap objects, and operators dispatch according to object behavior. This section states the direct consequence for Python's type system.

Python is dynamic at the level of names. A name can first refer to an integer, later to a string, later to \texttt{None}, and the interpreter permits that rebinding without complaint. But Python is not loose at the level of operations. Once an expression executes, the actual objects involved must support that operation through their type's method table.

This is why the same name can legally move from \texttt{42} to \texttt{"42"}, but the expression \texttt{"42" + 1} still fails with \texttt{TypeError}. Python does not silently guess whether the programmer wanted numeric addition or string concatenation.

\begin{quote}
Names are flexible. Objects are typed. Operations are checked.
\end{quote}

This distinction explains reassignment freedom, the roles of \texttt{type()} and \texttt{isinstance()}, defined numeric promotion among \texttt{bool}/\texttt{int}/\texttt{float}, and why unrelated combinations such as text plus number raise errors instead of being silently coerced.
